\documentclass[12pt]{article}
\usepackage[T1]{fontenc}
\usepackage[utf8]{inputenc}
\usepackage{lmodern}
\usepackage{textcomp}
\usepackage{enumitem}
\usepackage{geometry}
\geometry{margin=1in}
\begin{document}
\begin{center}
Answer Key - History - K-12 Level Assessment 19 \
\small (Answers are ordered exactly as in the question paper.)
\end{center}

\section*{Multiple Choice}
\begin{enumerate}[label=\arabic*.]
\item \textbf{Answer:} C
\\ \textit{Options:} A) sunflower; B) squash; C) maize; D) knotweed
\\ \textit{Note:} Maize was not an important component of Hopewell subsistence because the Hopewell culture primarily relied on native plants such as wild seeds, nuts, and tubers, as well as hunting and fishing, before maize became a staple in later agricultural societies.
\item \textbf{Answer:} D
\\ \textit{Options:} A) many processes that can be observed act extremely slowly.; B) uniformitarianism can account for significant geological features.; C) erosion works incredibly slowly.; D) All of the above.
\\ \textit{Note:} Charles Lyell's argument for the Earth's extreme age is supported by his observations of geological processes, the layering of rock strata, and the fossil record, all of which indicate that natural events occur gradually over long periods, thus leading to the conclusion that the Earth itself must also be very old.
\item \textbf{Answer:} C
\\ \textit{Options:} A) his horses and guns frightened the Incas into submission.; B) he was such a brilliant military strategist.; C) the empire had never won the allegiance of the peoples it had conquered.; D) he arrived with a superior force of more than 80,000 soldiers.
\\ \textit{Note:} The Inca empire's failure to secure the loyalty of the diverse groups it conquered left many of those peoples disillusioned and willing to ally with Pizarro against their Inca rulers.
\item \textbf{Answer:} B
\\ \textit{Options:} A) Khipu; B) Mit'a; C) Split inheritance; D) Capacocha
\\ \textit{Note:} The correct answer is "Mit'a" because it refers to the system of mandatory public service in the Incan Empire where subjects were required to work on state projects, including agriculture and infrastructure, in exchange for the right to use land.
\item \textbf{Answer:} C
\\ \textit{Options:} A) the use of seeds as a natural insecticide in Altamira cave in southwestern Europe; B) plant materials such as flowers in Aurignacian burials in southern France; C) plant material lining the floor of Sibudu Cave in South Africa; D) the use of wood in the Gravettian tradition to make spear-throwers and canoes
\\ \textit{Note:} The plant material lining the floor of Sibudu Cave in South Africa serves as evidence of new uses for plants during the Middle Stone Age, indicating that early humans utilized vegetation for purposes such as bedding or insulation, reflecting an advanced understanding of their environment.
\end{enumerate}

\section*{True / False}
\begin{enumerate}[label=\arabic*.]
\setcounter{enumi}{5}
\item \textbf{Answer:} False
\\ \textit{Options:} A) True; B) False
\\ \textit{Note:} The Dorset culture is actually known for its use of specialized tools such as ground slate and not fluted stone blades, and they primarily hunted smaller marine mammals rather than large terrestrial mammals like extinct elephants.
\item \textbf{Answer:} True
\\ \textit{Options:} A) True; B) False
\\ \textit{Note:} The statement is true because archaeological evidence indicates that the Yang-shao culture, known for its farming and pottery, evolved into the Lung-Shan culture, characterized by advancements in technology and social organization, around 5,000 years before present.
\item \textbf{Answer:} True
\\ \textit{Options:} A) True; B) False
\\ \textit{Note:} Archaeological findings in Israel, including ancient tools and remains of wild cereals, provide evidence that humans were harvesting and consuming wild grains in the region as far back as 20,000 years ago.
\item \textbf{Answer:} False
\\ \textit{Options:} A) True; B) False
\\ \textit{Note:} The statement is false because Lewis Henry Morgan proposed a theory of cultural evolution that included stages of Savagery, Barbarism, and Civilization, but not Collapse.
\item \textbf{Answer:} True
\\ \textit{Options:} A) True; B) False
\\ \textit{Note:} Uniformitarianism posits that the same geological processes observed in the present, such as weathering and erosion, have been shaping the Earth's surface over long periods, thus supporting the idea that these processes are consistent over time.
\end{enumerate}

\section*{Long Form Response}
\begin{enumerate}[label=\arabic*.]
\setcounter{enumi}{10}
\item \textbf{Answer:} Recent discoveries at the Cuncaicha rockshelter and Pucuncho quarry are significant in understanding the habitats inhabited by Paleoindians in the New World. These sites reveal that Paleoindians lived in a wide diversity of environments, including areas at some of the highest elevations in the Americas. The evidence found at these locations, such as tools and remnants of animal remains, suggests that these early inhabitants were adaptable and skilled at surviving in challenging landscapes. By studying these findings, researchers can better appreciate how Paleoindians interacted with their environments and the strategies they employed to thrive in various habitats. This information highlights the complexity of early human life in the Americas and broadens our understanding of their migration and settlement patterns.
\\ \textit{Note:} correct because it highlights the adaptability of Paleoindians to diverse and challenging environments, as evidenced by the tools and animal remains found at Cuncaicha rockshelter and Pucuncho quarry, thereby enhancing our understanding of their habitats in the New World.
\item \textbf{Answer:} The analysis of animal teeth from a nearby village reveals that the people who built Stonehenge brought animals to the site from outside their immediate region. This suggests that they engaged in trade or long-distance movement, indicating a network of connections with other communities. The variety of animal remains found at the site points to diverse agricultural practices and the importance of livestock in their society. Additionally, this evidence helps historians understand the social and economic interactions of the people during that time, highlighting how they might have traveled and exchanged resources. Overall, the study of these animal teeth provides valuable insights into the lifestyle and practices of the builders of Stonehenge.
\\ \textit{Note:} correct because it connects the analysis of animal teeth to evidence of trade and long-distance movement, illustrating the social and agricultural practices of the people who built Stonehenge.
\end{enumerate}

\vfill
\noindent\textit{End of Answer Key}
\end{document}